\SourceLeaf{13}
\[
\frac{\mathrm{d}s'.\cos MQ'}{r'r'}+\frac{\mathrm{d}s''.\cos MQ''}{r''r''}+\frac{\mathrm{d}s'''.\cos MQ'''}{r'''r'''}+\text{ etc.}=0
\]
im Fall II dagegen (weil die Zahl der Gleichungen ungerade ist)
\[
\frac{\mathrm{d}s'.\cos MQ'}{r'r'}+\frac{\mathrm{d}s''.\cos MQ''}{r''r''}+\frac{\mathrm{d}s'''.\cos MQ'''}{r'''r'''}+\text{ etc.}=-\mathrm{d}\Sigma
\]
\par
Behandeln wir alle Elemente $\mathrm{d}\Sigma$ auf dieselbe Weise und summieren, so erhalten wir links offenbar das über die ganze Oberfläche des Körpers erstreckte Integral $\int\frac{\mathrm{d}s.\cos MQ}{rr}$. Rechts erhalten wir im ersten Fall $0$ und im zweiten den gesamten Flächeninhalt der Kugeloberfläche mit Radius $=1$ mit negativem Vorzeichen, also $-4\pi$, wobei $\pi$ den halben Umfang eines Kreises mit Radius $=1$ bezeichnet.

Der Fall, in dem $M$ auf der Oberfläche des Körpers selbst liegt, ist gesondert zu behandeln. Man denke sich die Tangentialebene an die Oberfläche des Körpers im Punkt $M$. Sie teilt die Kugeloberfläche in zwei gleiche Halbkugelflächen. Die eine liegt auf derselben Seite der Ebene wie der feste Teil des Körpers bei $M$, die andere auf der gegenüberliegenden Seite. Für alle Elemente $\mathrm{d}\Sigma$ auf der ersten Halbkugelfläche ist der Punkt $M$ als innerer Punkt anzusehen, für die übrigen als äußerer. Daraus folgt, dass die Summe aller Ausdrücke
\[
\frac{\mathrm{d}s'.\cos MQ'}{r'r'}+\frac{\mathrm{d}s''.\cos MQ''}{r''r''}+\frac{\mathrm{d}s'''.\cos MQ'''}{r'''r'''}+\text{ etc.}
\]
nur den halben Flächeninhalt der Kugeloberfläche mit negativem Vorzeichen ergibt. Damit zeigen wir
\par\medskip\Needspace{4\baselineskip}
\begin{center}VIERTER LEHRSATZ.\end{center}
\par
\textit{Das über die ganze Oberfläche des Körpers erstreckte Integral $\int\frac{\mathrm{d}s.\cos MQ}{rr}$ ist $=0$, $=-2\pi$ oder $=-4\pi$, je nachdem, ob $M$ außerhalb des Körpers, auf seiner Oberfläche oder in seinem Inneren liegt.}
\par
{\rightskip=0pt plus 1em\noindent\hspace*{\parindent}Mit derselben Überlegung lässt sich ferner zeigen, dass das Integral $\int\frac{P\mathrm{d}s.\cos MQ}{rr}$ im ersten Fall allgemein verschwindet, wenn $P$ eine beliebige rationale Funktion der Größen $\cos MX$, $\cos MY$, $\cos MZ$ bezeichnet.\par}

\Article{7}
\par
Die Volumina des Kegelraums vom Scheitel bis zu den jeweiligen Punkten $P',P'',P'''$ usw. sind
\[
=\tfrac13r'\mathrm{d}\sigma',\quad\tfrac13r''\mathrm{d}\sigma'',\quad\tfrac13r'''\mathrm{d}\sigma'''\text{ etc.}
\]
oder
\[
=\pm\tfrac13r'\mathrm{d}s'.\cos MQ',\quad\mp\tfrac13r''\mathrm{d}s''.\cos MQ'',\quad\pm\tfrac13r'''\mathrm{d}s'''.\cos MQ'''\text{ etc.}
\]
